\documentclass[10pt,a4paper]{amsart}
\usepackage[margin=19mm]{geometry}
\usepackage{amsmath,amssymb,mathtools}
\usepackage[scr=rsfs,cal=euler]{mathalfa}
\usepackage{xfrac}
\usepackage[unicode,hidelinks,bookmarks=false]{hyperref}
\newcommand{\ie}{i.e.\ }
\newcommand{\cf}{cf.\ }
\newcommand{\resp}{respectively\ }
\newcommand{\Subsection}{\subsection}
\providecommand{\nobreakdash}{\nobreak\hbox{-}}
\setlength{\emergencystretch}{2em}
\multlinegap0pt
\usepackage{bm}
\usepackage{bbm}
\usepackage{dsfont}
\usepackage{xspace}
\usepackage{cancel}
\usepackage{mathtools}
\usepackage{amsthm}
\makeatletter
\@ifundefined{theorem}{\newtheorem{theorem}{Theorem}}{}
\@ifundefined{proposition}{\newtheorem{proposition}[theorem]{Proposition}}{}
\makeatother
\title[A note on finiteness of Tate cohomology groups]{A note on finiteness of Tate cohomology groups}
\author{Sabyasachi Dhar, Santosh Nadimpalli}
\date{Source: arXiv:2509.13297v1 (CC BY 4.0)}
\begin{document}
\maketitle

\noindent\emph{Source and licence.} A note on finiteness of Tate cohomology groups. Version arXiv:2509.13297v1 (CC BY 4.0), distributed under the Creative Commons Attribution 4.0 International licence, as recorded in the arXiv licence metadata for this exact version and shown on the abstract page https://arxiv.org/abs/2509.13297v1. Full rights notice: licenses/arxiv-2509.13297-CC-BY-4.0.txt.\par\smallskip

\noindent\textit{Editorial scope.} Counted unit: the Proposition whose source label is \texttt{prop} in the article (source0.tex lines 157--162, source label \texttt{prop}), delivered together with the article's own complete original proof of it (source0.tex lines 163--241, one \texttt{proof} environment of 79 lines). No mathematics is written, repaired or re-derived: statement and proof are the article's own text line for line. Required context retained uncounted: none. Every definition, display and label of those lines is retained unchanged. Omitted from this package and not redistributed: 1--156, 242--678. Nothing inside a retained excerpt is omitted or altered: every retained body is delivered exactly as the article's lines are written, so no omission receipt of any kind accompanies this package. Citations: the retained excerpts carry no citation command at all, so no bibliography entry of the article is transcribed and no reference is redistributed. Reference adaptation: none was needed and none was made; every reference inside the delivered unit resolves inside it, so no unresolved reference or citation is produced. Printed slips kept literal: no printed word, symbol, hypothesis or number of the counted statement or of its proof was altered. Layout and class: the article's own class is \texttt{amsart} with packages including lipsum, a4wide, amssymb, amsmath, amsthm, tikz, amscd, hyperref, enumitem; the wrapper is the lane's pinned \texttt{amsart} editorial wrapper at 10pt on A4, which restates the theorem-like environments the retained text needs and adds the article's own macro definitions that the retained text needs (none), with extra wrapper packages (bm, bbm, dsfont, xspace, cancel, mathtools). The delivered numbering is the unit's own. Assets: no retained span contains a figure environment, a graphics-inclusion command, a PDF-page inclusion, a table or a TikZ picture; no image file is copied into this package, the package declares no asset dependency and no image file is redistributed with it. Licence: version arXiv:2509.13297v1 of this article is distributed under the Creative Commons Attribution 4.0 International licence, as recorded in the arXiv licence metadata for this exact version and shown on the abstract page https://arxiv.org/abs/2509.13297v1. A full rights notice is delivered at licenses/arxiv-2509.13297-CC-BY-4.0.txt. Characters: every retained excerpt is pure ASCII, so the delivered engine, XeLaTeX with its default Latin Modern text font, reports no missing character.
\begin{proposition}\label{prop}
  Let $\Pi$ be a smooth $\ell$-modular representation of $G(F)$ such
  that $\Pi$ extends as a representation of $G(F)\rtimes \Gamma$. Then
  the Tate cohomology space $\widehat{H}^i(\Gamma, \Pi(U(F)))$ is
  trivial, for $i\in \{0,1\}$.
\end{proposition}

\begin{proof}
  The above proposition can be derived from Glauberman's
  correspondence applied to the $p$-group $U_0/U_1$, where $U_0$ and
  $U_1$ are $\sigma$-stable compact open subgroups of $U(F)$ such that
  $U_1$ is a normal subgroup of $U_0$. Since $U(F)$ does not have
  $\sigma$-fixed points, and from the fact that $U_0$ is a pro-$p$
  group, we get that $(U_0/U_1)^\sigma$ is a trivial group.  If
  $(\rho, V)$ is an $\ell$-modular representation of
  $(U_0/U_1)\rtimes \Gamma$, then we see that the Tate cohomology
  $\widehat{H}^i(\Gamma, \rho)$ decomposes into Tate cohomology of
  $\sigma$-stable sub-representations of $\rho$. But the only
  $\sigma$-stable sub-representations of $U_0/U_1$ are trivial
  representations by Glauberman's correspondence. We can apply these
  arguments to $\Pi(U(F))$ for $\sigma$-stable representations spanned
  by vectors which are annihilated by the augmentation ideal of some
  $U_0/U_1$ as above. We choose to give a direct argument below
  without appealing to Glauberman's correspondence.

Let us first consider the case $i=0$.
Let $w\in \Pi(U(F))^\sigma$ and let $U_0$ be a $\sigma$-stable
 compact open subgroup of $U(F)$
 such that 
 \begin{equation}\label{argu}
\int_{U_0}\Pi(u)w\,du=0,
 \end{equation}
 where $du$ is a Haar measure on $U_0$. Let $U_1$ be a $\sigma$-stable
 open normal subgroup of $U_0$ such that $w\in \Pi^{U_1}$. Let $H$ be
 the quotient group $U_0/U_1$. Then we have
$$ \int_{U_0}\Pi(u)w\,du = \sum_{h\in H} \Pi(h)w = \sum_{h\in H^\sigma} \Pi(h)w +
\sum_{h\in H\setminus H^\sigma} \Pi(h)w. $$ Since the $\Gamma$-action
on $H\setminus H^\sigma$ is free, there exists a subset
$X\subset H\setminus H^\sigma$ such that
$H\setminus H^\sigma = \coprod_{i=0}^{\ell-1}\sigma^iX$, and we get
the following equality
$$
\sum_{h\in H} \Pi(h)w = \sum_{h\in H^\sigma} \Pi(h)w + {\rm
  Nr}\big(\sum_{x\in X}\Pi(x)w\big).
$$
Note that $H^\sigma = \{e\}$ as $U_i$ are $p$-groups and $\sigma$ is
an element of order $\ell$. Then it follows from the above equation
that
$$ \sum_{h\in H} \Pi(h)w = w $$
in $\widehat{H}^0(\Gamma, \Pi(U(F)))$. Now, using the equation
(\ref{argu}), we get that $w=0$ in $\widehat{H}^0(\Gamma,\Pi(U(F)))$.

We now consider the case where $i=1$. 
Let $v\in {\rm Ker}\big({\rm Nr}:\Pi(U(F))\rightarrow \Pi(U(F))\big)$. There 
exists a $\sigma$-stable compact open subgroup $N_0$ of $U(F)$ such that
$$ \int_{N_0} \Pi(n)v\,dn = 0, $$
where $dn$ is a Haar measure on $N_0$. Let $N_1$ be a $\sigma$-stable
open normal subgroup of $N_0$ such that $v\in \Pi^{N_1}$. If
$\mathcal{N}$ denotes the quotient group $N_0/N_1$, then we have
$$ \int_{N_0} \Pi(n)v\,dn = \sum_{n\in \mathcal{N}} \Pi(n)v 
= \sum_{n\in \mathcal{N}^\sigma}\Pi(n)v + \sum_{n\in \mathcal{N}\setminus 
\mathcal{N}^\sigma}\Pi(n)v. $$
Note that $\mathcal{N}^\sigma = \{e\}$ as $N_i$ are $p$-groups and $\sigma$ 
is an element of order $\ell$, and there exists a subset $X\subseteq 
\mathcal{N}\setminus \mathcal{N}^\sigma$ such that $\mathcal{N}\setminus 
\mathcal{N}^\sigma = \coprod_{i=0}^{\ell-1}\sigma^iX$. Then the above 
identity becomes
\begin{equation}\label{argu_3}
v \,\,+ \,\,\sum_{i=0}^{\ell-1}\sum_{x\in X}\Pi(\sigma^i(x))v = 0.
\end{equation}
Set $w = \sum_{i=0}^{\ell-1}\sum_{x\in X}\Pi(\sigma^i(x))v$. We will show 
that $w\in {\rm Im}(1-\sigma)$. Using the fact that ${\rm Nr}(v) = 0$, we have
\begin{align*}
w 
&= \sum_{x\in X}\Pi(x)v \,\,+\,\, \sum_{i=1}^{\ell-1}\sum_{x\in X} \Pi(\sigma^i(x))v\\
&= -\sum_{i=1}^{\ell-1}\sum_{x\in X}
\Pi(x)\sigma^{i}v \,\,+\,\, \sum_{i=1}^{\ell-1}\sum_{x\in X} \Pi(\sigma^i(x)) v\\
&= -\sum_{i=1}^{\ell-1}\sigma^i\sum_{x\in X} \Pi(\sigma^{\ell-i}(x))v 
\,\,+\,\, \sum_{i=1}^{\ell-1}\sum_{x\in X}
\Pi(\sigma^{\ell-i}(x)) v\\
&= \sum_{i=1}^{\ell-1}(1-\sigma^i)\sum_{x\in X}\Pi^{\sigma^{\ell-i}}(x)v.
\end{align*}
Thus, we get that $w = (1-\sigma)w_1$ for some $w_1\in \Pi(U(F))$, and
hence it follows from (\ref{argu_3}) that $v = 0$ in
$\widehat{H}^1(\Gamma, \Pi(U(F))$.
\end{proof}

\end{document}
